% GENERATED FILE. Edit canonical lessons, then run npm run generate:books.
% canonical-source-hash: fnv1a64:354fa974c9e1ae38
% canonical-lessons: MW-C221-meh, MW-C221-barasno, MW-C221-bijli, MW-C221-gajno, MW-C221-jhari

\chapter{Rain to Come}
\label{ch:mw-rain-to-come}
\hlchaptermodality{221}

\begin{chapteropening}
\textbf{By the end of this chapter:} \emph{I can talk about the rain, thunder and lightning to come.}

Complete the last lesson of chapter 221: I can talk about the rain, thunder and lightning to come.
\end{chapteropening}

\section[\texorpdfstring{\mw{मेह} (meh)}{मेह (meh)}]{\mw{मेह} (meh) — the rains; rain}

\label{lesson:MW-C221-meh}



\begin{quote}
Before the new one: say the Marwadi for to change, then the Marwadi for to find out.
\end{quote}

\subsection*{You'll want to know: \mw{मेह}}
\textbf{\mw{मेह}} — \emph{meh} — \textquotedblleft{}the rains; rain\textquotedblright{}.

\begin{quote}
\mw{सावण} \mw{मांय} \mw{मेह} \mw{आवै} \mw{है।}
\end{quote}

\begin{quote}
\emph{sāvaṇ mā̃y meh āvai hai.}
\end{quote}

\begin{quote}
The rains come in the monsoon month.
\end{quote}

A masculine noun.

\begin{cousinweb}
\textbf{\mw{मेह}} (\emph{meh}) is inherited from Sanskrit \emph{megha}, \textquotedblleft{}a cloud, rain\textquotedblright{}; CDIAL lists Old Marwari \emph{meha}, \textquotedblleft{}rain\textquotedblright{} (CDIAL 10302).
\end{cousinweb}

\begin{grammarlens}[title={What you've built}]
One more word for saying what comes next.
\end{grammarlens}

\subsection*{Your turn}
\begin{itemize}
\raggedright
  \item \emph{Say it:} \emph{meh}
  \item \emph{Say it:} \emph{meh}, once more
  \item \emph{Say it:} \emph{sāvaṇ mā̃y meh āvai hai.} — The rains come in the monsoon month.
  \item \emph{Recall:} say the Marwadi for to change, then the Marwadi for to find out, then say \emph{meh} again
\end{itemize}

\begin{tcolorbox}[breakable,colback=teal!4,colframe=teal!35!black,title={Before you move on}]
What does \mw{मेह} mean? (\textquotedblleft{}The rains; rain\textquotedblright{}.) Say it once more.
\end{tcolorbox}
\section[\texorpdfstring{\mw{बरसणो} (barasṇo)}{बरसणो (barasṇo)}]{\mw{बरसणो} (barasṇo) — to rain; to pour down}

\label{lesson:MW-C221-barasno}



\begin{quote}
Before the new one: say the Marwadi for to find out, then the Marwadi for the rains.
\end{quote}

\subsection*{You'll want to know: \mw{बरसणो}}
\textbf{\mw{बरसणो}} — \emph{barasṇo} — \textquotedblleft{}to rain; to pour down\textquotedblright{}.

\begin{quote}
\mw{सावण} \mw{मांय} \mw{मेह} \mw{घणो} \mw{बरसै} \mw{है।}
\end{quote}

\begin{quote}
\emph{sāvaṇ mā̃y meh ghaṇo barsai hai.}
\end{quote}

\begin{quote}
In the monsoon month it rains hard.
\end{quote}

Past (it …): \mw{वो} \mw{मेह} \mw{बरस्यो}, \emph{vo meh barsyo}.

\begin{cousinweb}
\textbf{\mw{बरसणो}} (\emph{barasṇo}) is inherited from Sanskrit \emph{varṣati}, \textquotedblleft{}rains\textquotedblright{}; CDIAL lists Hindi \emph{barasnā}, \textquotedblleft{}to rain\textquotedblright{}. CDIAL gives no Marwari form in this entry (CDIAL 11394). \textbf{\mw{बरसात}}, rain, belongs to the same family.
\end{cousinweb}

\begin{grammarlens}[title={What you've built}]
One more word for saying what comes next.
\end{grammarlens}

\subsection*{Your turn}
\begin{itemize}
\raggedright
  \item \emph{Say it:} \emph{barasṇo}
  \item \emph{Say it:} \emph{barasṇo}, once more
  \item \emph{Say it:} \emph{sāvaṇ mā̃y meh ghaṇo barsai hai.} — In the monsoon month it rains hard.
  \item \emph{Recall:} say the Marwadi for to find out, then the Marwadi for the rains, then say \emph{barasṇo} again
\end{itemize}

\begin{tcolorbox}[breakable,colback=teal!4,colframe=teal!35!black,title={Before you move on}]
What does \mw{बरसणो} mean? (\textquotedblleft{}To rain; to pour down\textquotedblright{}.) Say it once more.
\end{tcolorbox}
\section[\texorpdfstring{\mw{बीजळी} (bījḷī)}{बीजळी (bījḷī)}]{\mw{बीजळी} (bījḷī) — lightning}

\label{lesson:MW-C221-bijli}



\begin{quote}
Before the new one: say the Marwadi for the rains, then the Marwadi for to rain.
\end{quote}

\subsection*{You'll want to know: \mw{बीजळी}}
\textbf{\mw{बीजळी}} — \emph{bījḷī} — \textquotedblleft{}lightning\textquotedblright{}.

\begin{quote}
\mw{बादल} \mw{मांय} \mw{बीजळी} \mw{है।}
\end{quote}

\begin{quote}
\emph{bādal mā̃y bījḷī hai.}
\end{quote}

\begin{quote}
There's lightning in the clouds.
\end{quote}

A feminine noun.

\begin{cousinweb}
\textbf{\mw{बीजळी}} (\emph{bījḷī}) is inherited from Sanskrit \emph{vidyullatā}, \textquotedblleft{}forked lightning\textquotedblright{}, built from \emph{vidyut}, \textquotedblleft{}lightning\textquotedblright{}, and \emph{latā}, \textquotedblleft{}a creeper\textquotedblright{}; CDIAL lists Old Marwari \emph{bījalī} (CDIAL 11745).
\end{cousinweb}

\begin{grammarlens}[title={What you've built}]
One more word for saying what comes next.
\end{grammarlens}

\subsection*{Your turn}
\begin{itemize}
\raggedright
  \item \emph{Say it:} \emph{bījḷī}
  \item \emph{Say it:} \emph{bījḷī}, once more
  \item \emph{Say it:} \emph{bādal mā̃y bījḷī hai.} — There's lightning in the clouds.
  \item \emph{Recall:} say the Marwadi for the rains, then the Marwadi for to rain, then say \emph{bījḷī} again
\end{itemize}

\begin{tcolorbox}[breakable,colback=teal!4,colframe=teal!35!black,title={Before you move on}]
What does \mw{बीजळी} mean? (\textquotedblleft{}Lightning\textquotedblright{}.) Say it once more.
\end{tcolorbox}
\section[\texorpdfstring{\mw{गाजणो} (gājṇo)}{गाजणो (gājṇo)}]{\mw{गाजणो} (gājṇo) — to thunder; to roar}

\label{lesson:MW-C221-gajno}



\begin{quote}
Before the new one: say the Marwadi for to rain, then the Marwadi for lightning.
\end{quote}

\subsection*{You'll want to know: \mw{गाजणो}}
\textbf{\mw{गाजणो}} — \emph{gājṇo} — \textquotedblleft{}to thunder; to roar\textquotedblright{}.

\begin{quote}
\mw{आज} \mw{रात} \mw{बादल} \mw{गाजै} \mw{है।}
\end{quote}

\begin{quote}
\emph{āj rāt bādal gājai hai.}
\end{quote}

\begin{quote}
Tonight the clouds are thundering.
\end{quote}

Past (it …): \mw{वो} \mw{बादल} \mw{गाज्यो}, \emph{vo bādal gājyo}.

\begin{cousinweb}
\textbf{\mw{गाजणो}} (\emph{gājṇo}) is inherited from Sanskrit \emph{garjati}, \textquotedblleft{}thunders, roars\textquotedblright{}; CDIAL lists the Old Marwari verb \emph{gājaï}, \textquotedblleft{}thunders\textquotedblright{} (CDIAL 4046).
\end{cousinweb}

\begin{grammarlens}[title={What you've built}]
One more word for saying what comes next.
\end{grammarlens}

\subsection*{Your turn}
\begin{itemize}
\raggedright
  \item \emph{Say it:} \emph{gājṇo}
  \item \emph{Say it:} \emph{gājṇo}, once more
  \item \emph{Say it:} \emph{āj rāt bādal gājai hai.} — Tonight the clouds are thundering.
  \item \emph{Recall:} say the Marwadi for to rain, then the Marwadi for lightning, then say \emph{gājṇo} again
\end{itemize}

\begin{tcolorbox}[breakable,colback=teal!4,colframe=teal!35!black,title={Before you move on}]
What does \mw{गाजणो} mean? (\textquotedblleft{}To thunder; to roar\textquotedblright{}.) Say it once more.
\end{tcolorbox}
\section[\texorpdfstring{\mw{झड़ी} (jhaṛī)}{झड़ी (jhaṛī)}]{\mw{झड़ी} (jhaṛī) — a spell of steady rain}

\label{lesson:MW-C221-jhari}



\begin{quote}
Before the new one: say the Marwadi for lightning, then the Marwadi for to thunder.
\end{quote}

\subsection*{You'll want to know: \mw{झड़ी}}
\textbf{\mw{झड़ी}} — \emph{jhaṛī} — \textquotedblleft{}a spell of steady rain\textquotedblright{}.

\begin{quote}
\mw{सावण} \mw{मांय} \mw{झड़ी} \mw{लागै} \mw{है।}
\end{quote}

\begin{quote}
\emph{sāvaṇ mā̃y jhaṛī lāgai hai.}
\end{quote}

\begin{quote}
In the monsoon month the rain sets in for days.
\end{quote}

A feminine noun.

\begin{cousinweb}
\textbf{\mw{झड़ी}} (\emph{jhaṛī}) is inherited from a reconstructed \emph{*jhaḍī}, \textquotedblleft{}a rainstorm\textquotedblright{}; CDIAL lists Marwari \emph{jhaṛī}, \textquotedblleft{}continuous heavy rain\textquotedblright{}, and points to a Dravidian word for fine steady rain (CDIAL 5329).
\end{cousinweb}

\begin{grammarlens}[title={What you've built}]
One more word for saying what comes next.
\end{grammarlens}

\subsection*{Your turn}
\begin{itemize}
\raggedright
  \item \emph{Say it:} \emph{jhaṛī}
  \item \emph{Say it:} \emph{jhaṛī}, once more
  \item \emph{Say it:} \emph{sāvaṇ mā̃y jhaṛī lāgai hai.} — In the monsoon month the rain sets in for days.
  \item \emph{Recall:} say the Marwadi for lightning, then the Marwadi for to thunder, then say \emph{jhaṛī} again
\end{itemize}

\begin{tcolorbox}[breakable,colback=teal!4,colframe=teal!35!black,title={Before you move on}]
What does \mw{झड़ी} mean? (\textquotedblleft{}A spell of steady rain\textquotedblright{}.) Say it once more.
\end{tcolorbox}
